%--------------------------------------------------------------------------
% CÓDIGO FUENTE RELEVANTE DEL APLICATIVO
%--------------------------------------------------------------------------
% Archivo: Apendices/CodigoFuente.tex
% Autor: Juan Diego Arévalo Pidiache
% Fecha: Marzo 2026
% Versión: CONTEXTUALIZADA AL REPOSITORIO REAL
%--------------------------------------------------------------------------

En este apéndice se incluyen fragmentos seleccionados del código real del aplicativo CERMONT. Se priorizan archivos que evidencian la lógica de negocio, la seguridad, la persistencia y la operación offline del sistema actualmente desarrollado en \texttt{cermont\_aplicativo/}.

%--------------------------------------------------------------------------
\section*{Máquina de estados de órdenes - \texttt{lib/utils/work-order-fsm.ts}}
%--------------------------------------------------------------------------

\begin{lstlisting}[language=JavaScript,caption={Estados y transiciones válidas de órdenes de trabajo},label={lst:work-order-fsm}]
export type WorkOrderStatus = 
  | 'borrador'
  | 'propuesta_enviada'
  | 'aprobada_con_po'
  | 'en_planeacion'
  | 'planificada'
  | 'lista_para_ejecucion'
  | 'en_curso'
  | 'pausada'
  | 'ejecucion_completada'
  | 'informe_elaborado'
  | 'revision'
  | 'ses_aprobada'
  | 'completada'
  | 'cerrada'
  | 'facturada'
  | 'cancelada'
  | 'pagada'
  | 'rechazada';

export const WORK_ORDER_TRANSITIONS: Record<WorkOrderStatus, WorkOrderStatus[]> = {
  borrador: ['propuesta_enviada', 'cancelada'],
  propuesta_enviada: ['aprobada_con_po', 'rechazada', 'borrador'],
  aprobada_con_po: ['en_planeacion', 'cancelada'],
  en_planeacion: ['planificada', 'cancelada'],
  planificada: ['lista_para_ejecucion', 'en_planeacion'],
  lista_para_ejecucion: ['en_curso', 'planificada'],
  en_curso: ['pausada', 'ejecucion_completada'],
  pausada: ['en_curso', 'ejecucion_completada'],
  ejecucion_completada: ['informe_elaborado', 'en_curso'],
  informe_elaborado: ['revision', 'ejecucion_completada'],
  revision: ['ses_aprobada', 'informe_elaborado'],
  ses_aprobada: ['completada', 'revision'],
  completada: ['cerrada', 'facturada'],
  cerrada: ['facturada', 'pagada'],
  facturada: ['pagada', 'cerrada'],
  pagada: [],
  cancelada: [],
  rechazada: ['borrador'],
};
\end{lstlisting}

%--------------------------------------------------------------------------
\section*{Autorización por rol - \texttt{lib/auth/rbac.ts} y \texttt{lib/auth/roles.ts}}
%--------------------------------------------------------------------------

\begin{lstlisting}[language=JavaScript,caption={Control de acceso por rol para rutas y endpoints sensibles},label={lst:rbac}]
export const ADMIN_ROLES: readonly UserRole[] = ["gerente", "administrativo"];

export type UserRole =
  | "gerente"
  | "residente"
  | "hes"
  | "supervisor"
  | "operador"
  | "tecnico"
  | "administrativo"
  | "cliente";

export async function withRole(
  allowedRoles: readonly UserRole[],
): Promise<AuthenticatedSession | NextResponse> {
  const session = await getSession();

  if (!session?.user) {
    return NextResponse.json({ error: "No autorizado" }, { status: 401 });
  }

  if (!hasRole(session.user.role, allowedRoles)) {
    return NextResponse.json(
      { error: "No tiene permisos para esta accion." },
      { status: 403 },
    );
  }

  return session as AuthenticatedSession;
}
\end{lstlisting}

%--------------------------------------------------------------------------
\section*{Modelo de datos principal - \texttt{prisma/schema.prisma}}
%--------------------------------------------------------------------------

\begin{lstlisting}[language=JavaScript,caption={Modelo principal de órdenes y enumeración de estados en Prisma},label={lst:prisma-work-orders}]
model work_orders {
  id                 String         @id
  numero_ot          String         @unique @db.VarChar(50)
  cliente            String         @db.VarChar(255)
  descripcion        String
  ubicacion          String?        @db.VarChar(500)
  estado             order_status   @default(borrador)
  prioridad          order_priority @default(media)
  fecha_inicio       DateTime?      @db.Timestamptz(3)
  fecha_fin          DateTime?      @db.Timestamptz(3)
  metadata           Json?
  created_at         DateTime       @default(now()) @db.Timestamptz(3)
  updated_at         DateTime       @db.Timestamptz(3)
  administrative_cierre administrative_cierre?
  cost_tracking      cost_tracking?
  evidences          evidences[]
  inspections        inspections[]
  invoices           invoices[]
  technical_reports  technical_reports[]
}

enum order_status {
  borrador
  propuesta_enviada
  aprobada_con_po
  en_planeacion
  planificada
  lista_para_ejecucion
  en_curso
  pausada
  ejecucion_completada
  informe_elaborado
  revision
  ses_aprobada
  completada
  cerrada
  facturada
  cancelada
  pagada
  rechazada
}
\end{lstlisting}

%--------------------------------------------------------------------------
\section*{Seguridad y limitación de tasa - \texttt{lib/middleware/rate-limiting.ts}}
%--------------------------------------------------------------------------

\begin{lstlisting}[language=JavaScript,caption={Políticas de rate limiting para autenticación y APIs},label={lst:rate-limiting}]
export const rateLimits = {
  auth: { name: 'auth', limit: 10, windowMs: 15 * 60_000 },
  login: { name: 'login', limit: 5, windowMs: 15 * 60_000 },
  api: { name: 'api', limit: 100, windowMs: 60_000 },
  upload: { name: 'upload', limit: 10, windowMs: 60 * 60_000 },
  orders: { name: 'orders', limit: 20, windowMs: 60 * 60_000 },
  reports: { name: 'reports', limit: 5, windowMs: 60 * 60_000 },
  public: { name: 'public', limit: 100, windowMs: 60_000 },
};

export async function applyRateLimit(
  identifier: string,
  policy: RateLimitPolicy
): Promise<Response | null> {
  const result = await checkRateLimit(identifier, policy);

  if (!result.success) {
    return new Response(
      JSON.stringify({ error: 'Too many requests', retryAfter: result.retryAfter }),
      { status: 429, headers: { 'Content-Type': 'application/json', ...result.headers } }
    );
  }

  return null;
}
\end{lstlisting}

%--------------------------------------------------------------------------
\section*{Cola offline y sincronización diferida - \texttt{lib/pwa/offline-queue.ts}}
%--------------------------------------------------------------------------

\begin{lstlisting}[language=JavaScript,caption={Persistencia local de acciones críticas para operación offline},label={lst:offline-queue}]
export interface OfflineAction {
  id: string;
  type:
    | 'CREATE_EVIDENCE'
    | 'UPDATE_ORDER_STATUS'
    | 'ADD_COST'
    | 'UPDATE_ORDER'
    | 'CREATE_INSPECTION'
    | 'UPDATE_INSPECTION'
    | 'UPSERT_CLOSURE';
  payload: Record<string, unknown>;
  createdAt: number;
  retries: number;
  status?: 'pending' | 'failed';
}

const DB_NAME = 'CermontOfflineDB';
const STORE_NAME = 'pending_actions';

export async function queueOfflineAction(
  action: Omit<OfflineAction, 'id' | 'createdAt' | 'retries'>
): Promise<void> {
  // Guarda la accion en IndexedDB para reintento posterior
}
\end{lstlisting}

Los fragmentos anteriores muestran cuatro aspectos centrales del sistema real: el FSM que regula el flujo de órdenes, el RBAC que condiciona acceso por rol, el modelo documental que soporta el ciclo operativo-administrativo y la cola offline que permite continuidad operativa en campo.
  retryCount: number;
  maxRetries: number;
  createdAt: Date;
  updatedAt: Date;
  lastError?: string;
}
// ... implementacion de la clase SyncQueueService
\end{lstlisting}

%--------------------------------------------------------------------------
\section*{Frontend: Servicio de Sincronización Offline}
%--------------------------------------------------------------------------

\begin{lstlisting}[language=JavaScript,caption={Servicio de Sincronización (Frontend) - \texttt{sync-service.ts}},label={lst:sync-service-frontend}]
import { openDB, DBSchema, IDBPDatabase } from 'idb';

interface OfflineDB extends DBSchema {
  pendingActions: {
    key: string;
    value: {
      id: string;
      type: 'CREATE' | 'UPDATE' | 'DELETE';
      endpoint: string;
      data: any;
      timestamp: number;
      retries: number;
    };
  };
  cachedData: {
    key: string;
    value: {
      id: string;
      type: string;
      data: any;
      timestamp: number;
    };
    indexes: {
      'by-type': string;
    };
  };
  // ... otros stores
}

class SyncService {
  private db: IDBPDatabase<OfflineDB> | null = null;

  async init() {
    if (this.db) return this.db;

    this.db = await openDB<OfflineDB>('cermont-offline', 1, {
      upgrade(db) {
        if (!db.objectStoreNames.contains('pendingActions')) {
          db.createObjectStore('pendingActions', { keyPath: 'id' });
        }
        // ...
      },
    });
    return this.db;
  }
  // ... metodos de sincronizacion
}
\end{lstlisting}

%--------------------------------------------------------------------------
\section*{Middleware de Autenticación JWT}
%--------------------------------------------------------------------------

\begin{lstlisting}[language=JavaScript,caption={Middleware de Autenticación - \texttt{authenticate.ts}},label={lst:auth-middleware}]
import { Request, Response, NextFunction } from 'express';
import { tokenBlacklistRepository } from '../../infra/db/repositories/TokenBlacklistRepository.js';
import { logger } from '../utils/logger.js';
import { loginFailedTotal } from '../metrics/prometheus.js';
import { jwtService } from '../security/jwtService.js';

declare global {
  namespace Express {
    interface Request {
      user?: {
        userId: string;
        email: string;
        role: string;
        jti: string;
      };
    }
  }
}

export async function authenticate(
  req: Request,
  res: Response,
  next: NextFunction
): Promise<void> {
  try {
    const authHeader = req.headers.authorization;

    if (!authHeader || !authHeader.startsWith('Bearer ')) {
      loginFailedTotal.inc({ reason: 'missing_token' });
      res.status(401).json({
        type: 'https://httpstatuses.com/401',
        title: 'Unauthorized',
        status: 401,
        detail: 'Token de autenticación requerido',
      });
      return;
    }

    const token = authHeader.substring(7);
    // ... verificacion de token
  } catch (error) {
    // ... manejo de errores
  }
}
\end{lstlisting}

\vspace{1cm}

\textbf{Nota:} Los fragmentos anteriores corresponden a la versión del código utilizada en las pruebas de validación del sistema. El código fuente completo se encuentra disponible en el repositorio del proyecto.

%--------------------------------------------------------------------------
% FIN DE CÓDIGO FUENTE
%--------------------------------------------------------------------------









